%==============================================================================
%  Cover Letter -- Paper 2 (Reliability & Calibration of Quantum-Kernel IDS)
%  GHI CHU: thay International Journal of Network Management bang ten tap chi Q2/Q3 da chon truoc khi nop.
%  Compile tren Overleaf: pdfLaTeX.
%==============================================================================
\documentclass[11pt]{letter}
\usepackage[utf8]{inputenc}
\usepackage[T1]{fontenc}
\usepackage[margin=1in]{geometry}
\usepackage{newtxtext}                 % Times, dong bo voi bai bao
\usepackage[hidelinks]{hyperref}       % link mau den, khong to xanh

\signature{%
  Quan Tran Anh Vo\\
  (on behalf of all co-authors)\\
  University of Science and Technology -- The University of Danang\\
  Danang 550000, Vietnam\\
  \href{mailto:votrananhquan1111@gmail.com}{votrananhquan1111@gmail.com}}

\address{%
  Quan Tran Anh Vo\\
  University of Science and Technology\\
  The University of Danang\\
  Danang 550000, Vietnam}

\begin{document}

\begin{letter}{%
  The Editor-in-Chief\\
  \textit{International Journal of Network Management}}

\opening{Dear Editor-in-Chief,}

We are pleased to submit our manuscript entitled \textbf{``Are Quantum-Kernel
Intrusion Detectors Trustworthy? A Reliability and Calibration Benchmark of QSVM
Against Strong Tabular Learners on NSL-KDD''} for consideration as a regular paper
in \textit{International Journal of Network Management}.

Quantum Support Vector Machines (QSVM) with the ZZ-feature map have been proposed as
a near-term quantum-machine-learning candidate for network intrusion detection (IDS).
However, existing studies evaluate them almost exclusively on \emph{accuracy} and only
against classical SVM baselines. On tabular IDS data, gradient-boosted trees and random
forests usually dominate accuracy, so an accuracy-only comparison misses the question
that actually governs deployment: whether the model's probabilistic alerts are
\emph{trustworthy}.

To address this gap, we reframe the evaluation around \emph{reliability and
calibration}. We benchmark a hardware-constrained four-qubit QSVM against the strongest
tabular learners -- Random Forest and XGBoost -- alongside an SVM-RBF baseline, under a
single zero-leakage preprocessing pipeline and a single Platt-scaling protocol, so that
observed differences are attributable to the classifier rather than to confounding
experimental factors. Using a five-run protocol on NSL-KDD, we report Expected
Calibration Error (ECE), Brier score, and AUC-PR with Cohen's $d$ effect sizes.

Our central finding is a \emph{regime-specific reliability} result: on the rare-attack
classes (User-to-Root and Remote-to-Local, under 1\% of the data and the most
security-critical), the quantum kernel is the most trustworthy model, attaining the
lowest ECE and Brier score with a large effect size against the tree ensembles. It is
also best calibrated at a balanced operating point and in the low-data regime, while
remaining only competitive under heavy prior shift and temporal drift -- results we
report transparently, including where the quantum kernel does not lead. We further show
that ranking quality and calibration are distinct (tree ensembles rank well yet are
over-confident), and that Platt scaling is the appropriate recalibration tool for
margin-based quantum and SVM models but not for tree ensembles.

We believe this manuscript is well aligned with the scope of \textit{International Journal of Network Management}.
It provides the first calibration-centric reliability study of a quantum kernel for
intrusion detection, together with practical, deployment-oriented guidance on when such
a model can be trusted -- a topic of growing relevance to readers working on machine
learning for network security and on trustworthy AI.

We confirm that:

\begin{itemize}
\item This manuscript is original and has not been published previously.
\item The manuscript is not currently under consideration by any other journal or
      conference.
\item All authors have reviewed and approved the manuscript and agree with its
      submission to \textit{International Journal of Network Management}.
\item There are no conflicts of interest related to this work.
\item The only self-citation is to our companion manuscript, which is clearly disclosed
      as such in the reference list.
\end{itemize}

This manuscript was prepared in \LaTeX{}; as requested for initial submission, we
provide a single compiled PDF for review and will supply the full \LaTeX{} source
files upon revision.

The corresponding author is \textbf{Quan Tran Anh Vo}
(\href{mailto:votrananhquan1111@gmail.com}{votrananhquan1111@gmail.com}). We thank you
and the reviewers in advance for your time and consideration, and we look forward to
your response.

\closing{Sincerely,}

\end{letter}
\end{document}
